\documentclass[a4paper,11pt]{article}
\usepackage[utf8]{inputenc}
\usepackage[spanish]{babel}
\usepackage{geometry}
\usepackage{enumitem}
\usepackage{sectsty}
\usepackage{titlesec}
\usepackage{tikz}
\usepackage{array}

\geometry{
    a4paper,
    left=20mm,
    top=18mm,
    right=20mm,
    bottom=20mm
}

% Sin números de página
\pagestyle{empty}

% Estilo de secciones con números en círculos
\newcommand{\circled}[1]{%
    \tikz[baseline=(char.base)]{%
        \node[shape=circle,draw,inner sep=0.5pt,fill=black,text=white,font=\bfseries\footnotesize] (char) {#1};%
    }%
}

% Formato de secciones
\titleformat{\section}
{\normalfont\Large\bfseries}
{\circled{\thesection}}{0.5em}{}

\titlespacing*{\section}{0pt}{12pt}{8pt}

\titleformat{\subsection}
{\normalfont\normalsize\bfseries}
{}{0em}{}

\titlespacing*{\subsection}{8pt}{4pt}{4pt}

\begin{document}

% Encabezado del curso
\begin{center}
\small{Universidad de Granada}

\vspace{0.3cm}

\textbf{\large{Producción e Interacción Oral - Nivel 7}}

\vspace{0.4cm}

\small
\begin{tabular}{@{}ll@{}}
\textbf{Grupos:} & A (Aula 24) y B (Aula 08) \\
\textbf{Centro:} & Centro de Lenguas Modernas / Edificio A \\
\textbf{Profesor:} & Javier Benítez Lainez \\
\textbf{Tutorías:} & Martes 9:00-10:00 y 14:30-15:30 \\
& Sala de profesores / Email: benitezl@go.ugr.es
\end{tabular}
\end{center}

\vspace{0.5cm}
\section*{Guía de Pretéritos: Narración en Pasado}

\vspace{0.4cm}

\section{Sistema Completo de Tiempos Pasados en Español C1}

\vspace{0.2cm}

---

\vspace{0.2cm}

\section{INTRODUCCIÓN}

\vspace{0.2cm}

\subsection{Propósito}

Esta guía presenta el sistema completo de tiempos pasados en español para narración oral nivel C1.

\vspace{0.2cm}

\subsection{Importancia}

El dominio de los pretéritos es esencial para narrar experiencias, historias y eventos pasados con precisión y sofisticación.

\vspace{0.2cm}

---

\vspace{0.2cm}

\section{LOS CUATRO PRETÉRITOS PRINCIPALES}

\vspace{0.2cm}

\subsection{Pretérito Indefinido (Pretérito Perfecto Simple)}

\vspace{0.2cm}

\section{Usos principales}

\begin{itemize}

  \item Acciones completadas en un momento específico del pasado

  \item Eventos con principio y fin claros

  \item Secuencia de acciones

\end{itemize}

\vspace{0.2cm}

\section{Terminaciones regulares}

\vspace{0.2cm}

\textbf{-AR verbos:}

\begin{itemize}

  \item -é, -aste, -ó, -amos, -asteis, -aron

  \item hablar → hablé, hablaste, habló, hablamos, hablasteis, hablaron

\end{itemize}

\vspace{0.2cm}

\textbf{-ER/-IR verbos:}

\begin{itemize}

  \item -í, -iste, -ió, -imos, -isteis, -ieron

  \item comer → comí, comiste, comió, comimos, comisteis, comieron

  \item vivir → viví, viviste, vivió, vivimos, vivisteis, vivieron

\end{itemize}

\vspace{0.2cm}

\section{Verbos irregulares comunes}

\begin{itemize}

  \item ser/ir → fui, fuiste, fue, fuimos, fuisteis, fueron

  \item estar → estuve, estuviste, estuvo, estuvimos, estuvisteis, estuvieron

  \item tener → tuve, tuviste, tuvo, tuvimos, tuvisteis, tuvieron

  \item hacer → hice, hiciste, hizo, hicimos, hicisteis, hicieron

  \item decir → dije, dijiste, dijo, dijimos, dijisteis, dijeron

\end{itemize}

\vspace{0.2cm}

\section{Marcadores temporales}

\begin{itemize}

  \item ayer, anteayer

  \item el martes pasado, el mes pasado, el año pasado

  \item en 1990, en 2015

  \item hace dos días, hace tres semanas

  \item el 5 de mayo

\end{itemize}

\vspace{0.2cm}

\section{Ejemplos}

\begin{itemize}

  \item "Ayer \textbf{hablé} con María."

  \item "\textbf{Hicimos} un viaje el año pasado."

  \item "\textbf{Fui} al mercado esta mañana."

\end{itemize}

\vspace{0.2cm}

---

\vspace{0.2cm}

\subsection{Pretérito Imperfecto}

\vspace{0.2cm}

\section{Usos principales}

\begin{itemize}

  \item Descripciones en el pasado

  \item Hábitos y rutinas pasadas

  \item Acciones en progreso cuando ocurrió otra

  \item Contexto o "fondo" de la historia

\end{itemize}

\vspace{0.2cm}

\section{Terminaciones regulares}

\vspace{0.2cm}

\textbf{-AR verbos:}

\begin{itemize}

  \item -aba, -abas, -aba, -ábamos, -abais, -aban

  \item hablar → hablaba, hablabas, hablaba, hablábamos, hablabais, hablaban

\end{itemize}

\vspace{0.2cm}

\textbf{-ER/-IR verbos:}

\begin{itemize}

  \item -ía, -ías, -ía, -íamos, -íais, -ían

  \item comer → comía, comías, comía, comíamos, comíais, comían

  \item vivir → vivía, vivías, vivía, vivíamos, vivíais, vivían

\end{itemize}

\vspace{0.2cm}

\section{Verbos irregulares (solo 3)}

\begin{itemize}

  \item ser → era, eras, era, éramos, erais, eran

  \item ir → iba, ibas, iba, íbamos, ibais, iban

  \item ver → veía, veías, veía, veíamos, veíais, veían

\end{itemize}

\vspace{0.2cm}

\section{Marcadores temporales}

\begin{itemize}

  \item antes, anteriormente

  \item en el pasado

  \item durante mi infancia

  \item siempre, normalmente, frecuentemente

  \item todos los días, cada semana

\end{itemize}

\vspace{0.2cm}

\section{Ejemplos}

\begin{itemize}

  \item "\textbf{Iba} al colegio todos los días."

  \item "\textbf{Era} muy feliz cuando era niño."

  \item "\textbf{Vivíamos} en Madrid antes."

\end{itemize}

\vspace{0.2cm}

---

\vspace{0.2cm}

\subsection{Pretérito Perfecto Compuesto}

\vspace{0.2cm}

\section{Usos principales}

\begin{itemize}

  \item Acciones pasadas con conexión al presente

  \item Periodos de tiempo NO completados

  \item Experiencias de vida

\end{itemize}

\vspace{0.2cm}

\section{Formación}

haber (presente) + participio

\vspace{0.2cm}

\textbf{Conjugación de haber:}

\begin{itemize}

  \item he, has, ha, hemos, habéis, han

\end{itemize}

\vspace{0.2cm}

\textbf{Participios regulares:}

\begin{itemize}

  \item -AR → -ado (hablado, comprado)

  \item -ER/-IR → -ido (comido, vivido)

\end{itemize}

\vspace{0.2cm}

\textbf{Participios irregulares comunes:}

\begin{itemize}

  \item abrir → abierto

  \item cubrir → cubierto

  \item decir → dicho

  \item escribir → escrito

  \item hacer → hecho

  \item morir → muerto

  \item poner → puesto

  \item resolver → resuelto

  \item romper → roto

  \item ver → visto

  \item volver → vuelto

\end{itemize}

\vspace{0.2cm}

\section{Marcadores temporales}

\begin{itemize}

  \item hoy, esta semana, este mes, este año

  \item últimamente, recientemente

  \item nunca, jamás (con experiencia)

  \item ya

  \item en los últimos años

\end{itemize}

\vspace{0.2cm}

\section{Ejemplos}

\begin{itemize}

  \item "\textbf{He trabajado} mucho esta semana."

  \item "\textbf{No he comido} nada hoy."

  \item "\textbf{He visitado} España tres veces."

\end{itemize}

\vspace{0.2cm}

---

\vspace{0.2cm}

\subsection{Pretérito Pluscuamperfecto}

\vspace{0.2cm}

\section{Usos principales}

\begin{itemize}

  \item Acción que ocurrió ANTES de otra acción pasada

  \item Expresión del pasado del pasado

\end{itemize}

\vspace{0.2cm}

\section{Formación}

haber (imperfecto) + participio

\vspace{0.2cm}

\textbf{Conjugación de haber (imperfecto):}

\begin{itemize}

  \item había, habías, había, habíamos, habíais, habían

\end{itemize}

\vspace{0.2cm}

\section{Ejemplos}

\begin{itemize}

  \item "\textbf{Había terminado} cuando llegaste."

  \item "Ya \textbf{habíamos comido} cuando llamaron."

  \item "\textbf{Había estudiado} mucho para el examen."

\end{itemize}

\vspace{0.2cm}

---

\vspace{0.2cm}

\section{DIFERENCIAS CLAVE}

\vspace{0.2cm}

\subsection{Indefinido vs. Imperfecto}

\vspace{0.2cm}

| Indefinido | Imperfecto |

|------------|------------|

| Acción completada | Descripción o hábito |

| "Ayer \textbf{fui} al cine" | "\textbf{Iba} al cine todos los sábados" |

| Evento específico | Estado o situación general |

| "Llovió \textbf{ayer}" | "\textbf{Llovía} cuando salí" |

| Principio y fin claros | Duración indefinida o repetición |

\vspace{0.2cm}

\section{Regla mnemotécnica}

\textbf{Indefinido = FOTO} (captura un momento específico)

\textbf{Imperfecto = VIDEO} (muestra escena, contexto, duración)

\vspace{0.2cm}

---

\vspace{0.2cm}

\subsection{Perfecto Compuesto vs. Indefinido}

\vspace{0.2cm}

| Perfecto Compuesto | Indefinido |

|--------------------|------------|

| Tiempo NO completado | Tiempo completado |

| "\textbf{He hablado} con ella \textbf{esta semana}" | "\textbf{Hablé} con ella \textbf{la semana pasada}" |

| Conexión con presente | Sin conexión con presente |

| "No \textbf{he comido} \textbf{hoy}" | "No \textbf{comí} \textbf{ayer}" |

\vspace{0.2cm}

---

\vspace{0.2cm}

\subsection{Pluscuamperfecto vs. Otros tiempos}

\vspace{0.2cm}

| Pluscuamperfecto | Indefinido | Imperfecto |

|------------------|------------|------------|

| Antes de otra acción pasada | Acción pasada simple | Simultáneo con otra acción |

| "\textbf{Había llegado} cuando \textbf{llegaste}" | "\textbf{Llegué} ayer" | "\textbf{Llegaba} cuando \textbf{llamaste}" |

\vspace{0.2cm}

---

\vspace{0.2cm}

\section{USOS AVANZADOS NIVEL C1}

\vspace{0.2cm}

\subsection{Narración de historias completas}

\vspace{0.2cm}

\textbf{Estructura ideal:}

1. \textbf{Imperfecto} (establecer contexto): "Era una noche oscura y llovía..."

2. \textbf{Indefinido} (eventos principales): "De repente, \textbf{escuché} un ruido..."

3. \textbf{Pluscuamperfecto} (antecedentes): "No \textbf{había escuchado} nada similar antes..."

4. \textbf{Imperfecto} (descripciones durante eventos): "El sonido \textbf{venía} del sótano..."

\vspace{0.2cm}

\section{Ejemplo completo}

"E\textbf{ra} una mañana soleada de primavera. \textbf{Desperté} temprano y \textbf{me levanté}. Mientras \textbf{desayunaba}, \textbf{escuché} las noticias. Ya \textbf{había decidido} que ese día sería especial, así que \textbf{salí} de casa con una sonrisa. El sol \textbf{brillaba} y las aves \textbf{cantaban}. \textbf{Caminé} hacia el parque donde \textbf{solía} encontrarme con mis amigos."

\vspace{0.2cm}

---

\vspace{0.2cm}

\section{ERRORES COMUNES}

\vspace{0.2cm}

\subsection{Error 1: Usar indefinido para hábitos}

❌ "\textbf{Fuí} al colegio todos los días."

✅ "\textbf{Iba} al colegio todos los días."

\vspace{0.2cm}

\subsection{Error 2: Usar imperfecto para evento específico}

❌ "\textbf{Iba} a la fiesta ayer."

✅ "\textbf{Fuí} a la fiesta ayer."

\vspace{0.2cm}

\subsection{Error 3: Confundir perfecto e indefinido}

❌ "Ayer \textbf{he hablado} con Juan."

✅ "Ayer \textbf{hablé} con Juan."

\vspace{0.2cm}

\subsection{Error 4: No usar pluscuamperfecto}

❌ "\textbf{Comí} antes de que \textbf{llegaras}." (confuso)

✅ "\textbf{Había comido} antes de que \textbf{llegaras}." (claro)

\vspace{0.2cm}

---

\vspace{0.2cm}

\section{EJERCICIOS PRÁCTICOS}

\vspace{0.2cm}

\subsection{Ejercicio 1: Elige el tiempo correcto}

\vspace{0.2cm}

1. "(Ayer/Iba) ____ al mercado cada semana.

\textbf{Respuesta:} Iba (hábito)

\vspace{0.2cm}

2. "(Ayer/Fui) ____ al mercado ayer.

\textbf{Respuesta:} Fui (evento específico)

\vspace{0.2cm}

3. "(He comido/Comí) ____ hoy a las 2.

\textbf{Respuesta:} He comido (hoy = tiempo no completado)

\vspace{0.2cm}

4. "(Comí/He comido) ____ ayer a las 2.

\textbf{Respuesta:} Comí (ayer = tiempo completado)

\vspace{0.2cm}

---

\vspace{0.2cm}

\subsection{Ejercicio 2: Narración completa}

\vspace{0.2cm}

Completa la historia con los tiempos apropiados:

\vspace{0.2cm}

"______ (Ser/Hacía) mucho calor. ______ (Decidir/Iba) ir a la playa. ______ (Llegar/Llegué) a las 10. Ya ______ (ser/estuvo) mucha gente allí. ______ (Ponerme/Puse) el traje de baño y ______ (entrar/entré) al agua. ______ (Estar/Estará) refrescante. Ya ______ (nadar/nadé) antes ese verano, pero esa vez ______ (disfrutar/disfrutaba) especialmente."

\vspace{0.2cm}

\textbf{Respuestas:} Hacía, Decidí, Llegué, había/había, Puse, entré, Estaba, había nadado, disfruté

\vspace{0.2cm}

---

\vspace{0.2cm}

\section{ESTRATEGIAS DE USO EN NARRACIÓN ORAL}

\vspace{0.2cm}

\subsection{Para narrar una historia}

1. \textbf{Inicia con imperfecto} (contexto): "Cuando era joven..."

2. \textbf{Usa indefinido} para eventos principales: "Un día, decidí..."

3. \textbf{Intercala imperfecto} para descripciones: "El día estaba soleado..."

4. \textbf{Usa pluscuamperfecto} para antecedentes: "Ya había pensado en eso..."

\vspace{0.2cm}

\subsection{Para describir hábitos pasados}

\begin{itemize}

  \item Solía + infinitivo: "\textbf{Solía} jugar al tenis."

  \item Imperfecto: "\textbf{Jugaba} al tenis todos los días."

\end{itemize}

\vspace{0.2cm}

\subsection{Para secuencia de eventos}

\begin{itemize}

  \item Primero, después, luego, finalmente + indefinido:

  \item "\textbf{Primero hice} esto, \textbf{después fui} allí, \textbf{luego compré}..."

\end{itemize}

\vspace{0.2cm}

---

\vspace{0.2cm}

\textbf{Sesión 10: Pretéritos - Narración en Pasado}

\textbf{Nivel C1 - Español Oral Avanzado}

\vspace{0.2cm}

\end{document}
